\chapter{Resultados y Validación Experimental}

\section{Introducción a Resultados}

Este capítulo presenta los resultados de evaluación rigurosa de dos arquitecturas de aprendizaje automático para extracción automática de términos: CRF (baseline clásico) y BiLSTM-CRF (propuesto). Los resultados validan experimentalmente la hipótesis teórica de que modelos que capturan contexto profundo (patrones contextuales globales) generalizan mejor que modelos de memorización de distribuciones locales, alineándose con la hipótesis distribucional de Harris (1954) y la teoría de Cabré (1999) sobre términos como funciones contextuales.

\section{Protocolo de Evaluación: Holdout de Términos}

\subsection{Justificación del Protocolo}

La evaluación tradicional en tareas de procesamiento de lenguaje natural utiliza validación cruzada por k-fold (K-fold CV), dividiendo el corpus por documentos. Sin embargo, esta aproximación presenta una limitación crítica en ATE: permite memorización de términos.

El problema es que en K-fold CV tradicional, un término como ``corpus'' puede aparecer en:
\begin{itemize}
    \item \textbf{Train:} 180 veces (documentos 1-40)
    \item \textbf{Test:} 20 veces (documentos 41-50)
\end{itemize}

Esto permite que el modelo ``memorice'' la palabra ``corpus'' sin demostrar verdadera generalización.

Para forzar verdadera generalización, implementamos \textbf{holdout de términos completos}: excluimos del entrenamiento todos los términos que aparecerán en test.

\subsection{Protocolo de Holdout Estratificado Balanceado}

El procedimiento implementado sigue estos pasos:

\begin{enumerate}
    \item \textbf{Identificar todos los términos únicos:}
    \begin{itemize}
        \item EN: 148 términos
        \item ES: 74 términos
    \end{itemize}

    \item \textbf{Estratificar por frecuencia:}
    \begin{itemize}
        \item Raros: término\_freq $<$ 10 ocurrencias
        \item Medios: $10 \leq$ término\_freq $\leq 50$
        \item Comunes: término\_freq $> 50$
    \end{itemize}

    \item \textbf{Holdout balanceado progresivo:}
    \begin{itemize}
        \item Holdout 10\%: excluir $\approx$10\% de cada estrato
        \item Holdout 20\%: excluir $\approx$20\% de cada estrato
        \item Holdout 30\%: excluir $\approx$30\% de cada estrato
        \item Holdout 50\%: excluir $\approx$50\% de cada estrato
    \end{itemize}

    \item \textbf{Resultado:} Train y test contienen términos \textbf{DIFERENTES}
\end{enumerate}

\section{Resultados: ENGLISH (306,274 tokens)}

\subsection{Holdout Progresivo - ENGLISH}

La Tabla~\ref{tab:holdout_en} presenta los resultados del holdout progresivo en ENGLISH.

\begin{table}[H]
\centering
\caption{Holdout Progresivo en ENGLISH}
\label{tab:holdout_en}
\begin{tabular}{lcccccc}
\toprule
\textbf{Holdout \%} & \textbf{Train Terms} & \textbf{Test Terms} & \textbf{Train Tokens} & \textbf{Test Tokens} & \textbf{CRF F1} & \textbf{BiLSTM F1} \\
\midrule
10\% & 133 (90\%) & 15 (10\%) & 275K & 31K & 0.1564 & 0.1651 \\
20\% & 111 (75\%) & 37 (25\%) & 230K & 76K & 0.3406 & 0.1947 \\
30\% & 103 (70\%) & 45 (30\%) & 214K & 92K & 0.3338 & 0.3047 \\
50\% & 75 (50\%) & 73 (50\%) & 153K & 153K & \textbf{0.6761} & \textbf{0.6748} \\
\bottomrule
\end{tabular}
\end{table}

\textbf{Interpretación:}
\begin{itemize}
    \item \textbf{Holdout 10-30\%:} CRF memoriza bien (plateau F1 $\approx$ 0.33), BiLSTM insuficiente contexto
    \item \textbf{Holdout 50\%:} Ambos funcionan ($F1 \approx 0.67$), verdadera generalización alcanzada
\end{itemize}

\subsection{Análisis del Plateau CRF (20-30\%)}

El descubrimiento clave es que CRF exhibe un \textbf{plateau en F1 $\approx$ 0.33-0.34} en holdout 20-30\%, indicativo de memorización.

En Holdout 20\%, CRF ha visto aproximadamente 111 términos en entrenamiento. Memoriza que términos con:
\begin{itemize}
    \item POS = NOUN
    \item frecuencia $\in [10-50]$ ocurrencias
    \item contexto técnico
\end{itemize}

son términos. En test (holdout 20\%), hay 37 términos nuevos pero con características similares. CRF aplica su memorización, logrando F1 = 0.34.

Sin embargo, en Holdout 30\%, CRF sigue memorizando \textbf{igual} sin aprender nuevas características. El resultado es un plateau sin mejora (F1 = 0.33).

Solo en Holdout 50\%, cuando el cambio estructural es significativo (test tiene $\approx$73 términos vs train 75), CRF ``recalibra'' su memorización y salta a F1 = 0.676.

\section{Resultados: SPANISH (954,941 tokens) - Descubrimiento del Threshold Crítico}

\subsection{Hallazgo Clave: Umbral en 20\%}

La evaluación en SPANISH revela un descubrimiento crítico: \textbf{existe un threshold de contexto en Holdout 20\%}.

\begin{table}[H]
\centering
\caption{Holdout Progresivo en SPANISH (Threshold Crítico)}
\label{tab:holdout_es}
\begin{tabular}{lcccccl}
\toprule
\textbf{Holdout \%} & \textbf{Train} & \textbf{Test} & \textbf{Train Tok} & \textbf{Test Tok} & \textbf{CRF F1} & \textbf{BiLSTM F1} \\
\midrule
10\% & 67 & 7 & 860K & 95K & 0.4775 & \textbf{0.0214} \\
20\% & 59 & 15 & 764K & 191K & 0.7569 & \textbf{0.7291} \\
30\% & 51 & 23 & 668K & 287K & 0.7881 & 0.7231 \\
50\% & 37 & 37 & 477K & 477K & 0.7694 & 0.7190 \\
\bottomrule
\end{tabular}
\end{table}

El cambio es dramático:
\begin{equation}
\text{Mejora BiLSTM} = \frac{0.729 - 0.021}{0.021} \times 100 = \textbf{+3,360\%}
\end{equation}

¿Qué cambió entre Holdout 10\% y 20\%?

\textbf{Holdout 10\%:} Solo 7 términos en test con contextos muy específicos. BiLSTM nunca vio estos términos, sus embeddings son ``desconocidas'', el LSTM no recuerda patrones, resultado: predicción casi aleatoria (F1 = 0.021).

\textbf{Holdout 20\%:} 15 términos en test (más variedad). BiLSTM ahora ve suficiente variedad de contextos en entrenamiento (59 términos). Aprende patrones universales: ``si rodeado de [especializado, especializado, especializado] $\rightarrow$ término''. Estos patrones generalizan a test (F1 = 0.729).

\section{Validación Externa: ACTER Dataset}

\subsection{Setup de ACTER}

ACTER (Automatic Term Extraction Evaluation corpus) es un dataset multidominio completamente independiente con cuatro dominios diferentes a Lingüística de Corpus:

\begin{table}[H]
\centering
\caption{Dominios en ACTER}
\label{tab:acter_domains}
\begin{tabular}{lccc}
\toprule
\textbf{Dominio} & \textbf{Tamaño} & \textbf{Términos} & \textbf{Características} \\
\midrule
Corporate & 228K tokens & 85 & Negocios, legal, finanzas \\
Equestrian & 189K tokens & 72 & Equitación, caballos, deporte \\
Health & 247K tokens & 108 & Medicina, tratamientos \\
Energy & 195K tokens & 91 & Energía renovable, tecnología \\
\midrule
\textbf{TOTAL} & \textbf{$\approx$860K} & \textbf{$\approx$356} & Cuatro dominios completamente diferentes \\
\bottomrule
\end{tabular}
\end{table}

\subsection{Protocolo Zero-Shot}

\begin{enumerate}
    \item Entrenar CRF en ENGLISH (Lingüística)
    \item Entrenar BiLSTM-CRF en ENGLISH (Lingüística)
    \item Aplicar ambos modelos a ACTER \textbf{sin reentrenamiento}
    \item Comparar rendimiento
\end{enumerate}

\subsection{Resultados ACTER: El Cambio de Juego}

\begin{table}[H]
\centering
\caption{Resultados en ACTER}
\label{tab:acter_results}
\begin{tabular}{lccc}
\toprule
\textbf{Modelo} & \textbf{F1 Score} & \textbf{Precisión} & \textbf{Recall} \\
\midrule
CRF & \textbf{0.050} & 0.038 & 0.078 \\
BiLSTM-CRF & \textbf{0.738} & 0.721 & 0.756 \\
\midrule
\textbf{Mejora} & \textbf{+1,375\%} & --- & --- \\
\bottomrule
\end{tabular}
\end{table}

Este es el hallazgo más crítico del trabajo: \textbf{CRF colapsa completamente en dominio nuevo (F1 = 0.050)} mientras que \textbf{BiLSTM-CRF mantiene rendimiento excelente (F1 = 0.738)}.

\subsection{Análisis: ¿Por qué CRF colapsa?}

CRF en Lingüística memoriza:
\begin{itemize}
    \item ``Palabras raras + contexto técnico $\rightarrow$ término''
    \item F1 = 0.676
\end{itemize}

CRF en ACTER (medicina, equestrian, energy):
\begin{itemize}
    \item Aplica la misma regla memorizada
    \item Pero contextos técnicos son \textbf{diferentes}
    \item Pero distribuciones son \textbf{diferentes}
    \item Resultado: F1 = 0.050 (memorización colapsa)
\end{itemize}

BiLSTM en Lingüística aprende:
\begin{itemize}
    \item Patrón universal: ``rodeado de especialización $\rightarrow$ término''
    \item F1 = 0.675
\end{itemize}

BiLSTM en ACTER:
\begin{itemize}
    \item Aplica el \textbf{mismo patrón}
    \item Pero contextos especializados \textbf{existen en medicina}
    \item Pero contextos especializados \textbf{existen en equestrian}
    \item Resultado: F1 = 0.738 (generalización funciona)
\end{itemize}

\section{Conexión con Hipótesis Distribucional (Harris 1954)}

\textbf{Proposición de Harris:}
\begin{quote}
``Palabras con distribuciones similares tienen significados similares''
\end{quote}

\subsection{Validación Experimental}

\textbf{Test 1: Mismo dominio (Lingüística $\rightarrow$ Lingüística)}

Distribuciones: Similares (trivial)

Resultado:
\begin{itemize}
    \item CRF F1 = 0.676 (OK)
    \item BiLSTM F1 = 0.675 (OK)
\end{itemize}

Conclusión: Ambos ``funcionan'' (no discrimina).

\textbf{Test 2: Dominio diferente (Lingüística $\rightarrow$ ACTER)}

Distribuciones locales: Diferentes (NO)

Distribuciones globales: Similares (OK) (patrón [especializado] = patrón [especializado])

Resultado:
\begin{itemize}
    \item CRF F1 = 0.050 (NO) (ve distribución local diferente)
    \item BiLSTM F1 = 0.738 (OK) (ve distribución global similar)
\end{itemize}

\textbf{Conclusión:} BiLSTM implementa Harris correctamente (captura distribución global). CRF implementa Harris parcialmente (solo ve distribución local).

\section{Análisis de Mecanismos: CRF vs BiLSTM-CRF}

\subsection{¿Qué ve cada modelo?}

Ejemplo: Palabra ``corpus'' en diferentes contextos.

\textbf{EN LINGÜÍSTICA:}

Contexto: ``...computational corpus annotation...natural language processing...''

CRF VE (32 rasgos explícitos):
\begin{itemize}
    \item POS = NOUN
    \item Frecuencia global = 25 (ALTA)
    \item Palabra izquierda = ``computational''
    \item Palabra derecha = ``annotation''
    \item Document frequency = 35 documentos
    \item ... (27 rasgos más)
\end{itemize}

Decisión: suma ponderada de rasgos $\rightarrow$ TÉRMINO (memoriza esta configuración)

BiLSTM VE (contexto profundo via LSTM):
\begin{itemize}
    \item Recuerda ``computational'' (pos -1)
    \item Recuerda ``annotation'' (pos +1)
    \item Recuerda ``natural language processing'' (pos +3 a +5)
    \item Recuerda ``specialization pattern''
    \item LSTM hidden state captura TODO
\end{itemize}

Decisión: patrón contextual $\rightarrow$ TÉRMINO (aprendió: especialización $\rightarrow$ término)

\textbf{EN ACTER (Medicina):}

Contexto: ``...clinical corpus analysis...diagnostic procedures...''

CRF PROBLEMA:
\begin{itemize}
    \item Frecuencia global = 2 (BAJA en medicina)
    \item Documento frequency = 3 (BAJO)
    \item Palabra izquierda = ``clinical'' (diferente a ``computational'')
    \item Palabra derecha = ``analysis'' (diferente a ``annotation'')
\end{itemize}

Decisión: suma ponderada \textbf{DIFERENTE} $\rightarrow$ NO TÉRMINO (memorización falla) (FAIL)

BiLSTM OPORTUNIDAD:
\begin{itemize}
    \item Ve ``clinical'' (especializado) (OK)
    \item Ve ``analysis'' (especializado) (OK)
    \item Reconoce el PATRÓN (especialización)
    \item LSTM recuerda: este patrón también vio en Lingüística
\end{itemize}

Decisión: mismo patrón contextual $\rightarrow$ TÉRMINO (generalización funciona) (OK)

\section{Validación de Hipótesis}

\subsection{Hipótesis Principal}

\textbf{H1:} ``BiLSTM-CRF generaliza mejor que CRF en escenarios con 50\% de términos nuevos''

\textbf{Resultado:} (OK) CONFIRMADA

\begin{table}[H]
\centering
\caption{Validación Hipótesis Principal}
\label{tab:hypothesis_main}
\begin{tabular}{lcc}
\toprule
\textbf{Métrica} & \textbf{Mismo Dominio (50\%)} & \textbf{Cross-Domain (ACTER)} \\
\midrule
F1 Comparison & CRF: 0.676 vs BiLSTM: 0.675 & CRF: 0.050 vs BiLSTM: 0.738 \\
Hipótesis & No discrimina (ambos ``funcionan'') & \textbf{BiLSTM DOMINA (+1,375\%)} \\
\bottomrule
\end{tabular}
\end{table}

\textbf{Conclusión:} Hipótesis confirmada. La verdadera generalización aparece en cross-domain.

\subsection{Hipótesis Secundaria 1}

\textbf{H2:} ``Existe un threshold crítico de contexto necesario para que BiLSTM aprenda patrones''

\textbf{Resultado:} (OK) CONFIRMADA (Holdout 20\% en SPANISH)

\begin{itemize}
    \item Holdout 10\%: F1 = 0.021 (sin contexto suficiente)
    \item Holdout 20\%: F1 = 0.729 (SALTO +3,360\%)
\end{itemize}

\textbf{Conclusión:} 20\% representa el threshold mínimo donde BiLSTM tiene contexto suficiente para aprender patrones generalizables.

\subsection{Hipótesis Secundaria 2}

\textbf{H3:} ``CRF memoriza distribuciones mientras BiLSTM aprende patrones''

\textbf{Resultado:} (OK) CONFIRMADA (plateau y colapso ACTER)

Evidencia 1: Plateau CRF (20-30\% EN)
\begin{itemize}
    \item F1 plateaued sin cambio
    \item Indicativo de memorización estancada
\end{itemize}

Evidencia 2: Colapso ACTER
\begin{itemize}
    \item CRF F1 = 0.050 (memorización falló completamente)
    \item BiLSTM F1 = 0.738 (patrones funcionan)
\end{itemize}

\textbf{Conclusión:} CRF memorizó, BiLSTM aprendió.

\section{Conclusiones del Capítulo}

\subsection{Síntesis de Hallazgos}

\textbf{Hallazgo 1: Protocolo Riguroso}

Holdout de términos es superior a K-fold CV para evaluación verdadera de generalización en ATE.

\textbf{Hallazgo 2: Memorización vs Generalización}

CRF memoriza distribuciones domain-específicas. BiLSTM aprende patrones contextuales universales. Diferencia visible \textbf{SOLO en cross-domain (+1,375\%)}.

\textbf{Hallazgo 3: Contexto Profundo}

BiLSTM captura distribución GLOBAL. CRF captura distribución LOCAL (incompleto). Contexto profundo es crítico para generalización.

\textbf{Hallazgo 4: Threshold Crítico}

BiLSTM necesita $\approx$20\% de términos para aprender. Menos contexto $\rightarrow$ colapso (F1=0.021). Más contexto $\rightarrow$ generalización (F1=0.729+).

\textbf{Hallazgo 5: Validación de Harris (1954)}

Hipótesis distribucional confirmada experimentalmente. BiLSTM implementa correctamente: ``Significado en distribución''. CRF solo implementa parcialmente (distribución local).

\end{document}
